\section{中美机器人分工：硬件、信仰和数据飞轮}

本章回到中美视角。谭杰认为中国硬件发展很快，宇树、智元、星海图等公司代表中国在本体、供应链和制造上的强项；美国/硅谷更强在长期投入、模型大脑、研究信仰和算力。理想状态下，中美应更好合作：美国的智能发展和中国的硬件制造互补。

\begin{figure}[H]
\centering
\includegraphics[width=0.94\textwidth]{china-us-robotics-split.png}
\caption{中美机器人分工：美国更强智能范式，中国更强硬件与供应链。自制概念图，依据 01:17:35--02:06:16 对谈内容整理。}
\end{figure}

\begin{knowledgebox}{读图：硬件强不自动等于数据飞轮强}
中国强在硬件、本体、供应链和制造；美国/硅谷强在长期研究、模型大脑、算力和愿意为远期愿景投入。机器人需要两者结合，但硬件必须先有实际用途，数据飞轮才会转起来。
\end{knowledgebox}

\subsection{硅谷信仰与国内短周期压力}

本节把技术路线转到资源环境。机器人需要长期数据和长期工程投入，因此不同地区的资本耐心、硬件能力和组织文化，会直接影响路线能不能跑出来。

谭杰说，硅谷愿意相信看起来 ambitious、短期没有结果的长期方向，可以投钱十年；国内更倾向短期落地、盈利和快速发展。机器人和大模型都需要烧很多钱、采很多数据、做长期试验。如果初期只给很少资源，很难验证 scaling。

\begin{warningbox}{短周期压力会伤害长期验证}
机器人许多路线不是“先给一点数据看效果，再决定是否投入”就能验证的。若初期数据量太小，模型甚至无法显现 scaling 规律，团队可能会错误地放弃长期正确方向。
\end{warningbox}

\subsection{Google 路线、Waymo 路线与特斯拉路线}

张小珺提出一个类比：Google 的机器人路线可能像 Waymo，强调大脑和系统；国内硬件公司可能像 Tesla，希望硬件部署带来数据。谭杰提醒，Tesla 的车本身有用，因此数据飞轮能转；当前很多机器人硬件还没有达到有用阈值，不能简单类比 Tesla。

\begin{knowledgebox}{数据飞轮成立的前提}
\begin{tabularx}{\textwidth}{p{0.24\textwidth}p{0.34\textwidth}X}
\toprule
前提 & 含义 & 机器人中的难点 \\
\midrule
产品有用 & 用户愿意真实使用，而不只是看 demo & 很多机器人还未达到日常可用阈值。 \\
场景高频 & 使用频率足以产生大量数据 & 家庭和开放场景长尾太多。 \\
反馈可记录 & 成功、失败和接管能被结构化记录 & 操作失败原因可能来自视觉、力控、规划或硬件。 \\
回流可训练 & 数据能进入训练和评测闭环 & 跨本体和隐私/安全都增加成本。 \\
\bottomrule
\end{tabularx}
\end{knowledgebox}

\subsection{本章小结}

中美机器人分工不是谁替代谁，而是硬件和智能的互补。真正的挑战是让硬件进入有用场景，形成数据飞轮，并用长期投入把模型、合成数据和真实评测接起来。

